\documentclass[tikz,border=10pt]{standalone}
\usepackage{tikz}
\usetikzlibrary{
    positioning,
    shapes.geometric,
    arrows.meta,
    calc,
    fit,
    backgrounds,
    decorations.pathreplacing,
    matrix
}

% Colors
\definecolor{receivercolor}{RGB}{255, 165, 0}
\definecolor{ringcolor}{RGB}{60, 179, 113}
\definecolor{hlcacolor}{RGB}{70, 130, 180}
\definecolor{lucacolor}{RGB}{205, 92, 92}
\definecolor{attncolor}{RGB}{255, 215, 0}
\definecolor{milcolor}{RGB}{180, 180, 220}

\begin{document}
\begin{tikzpicture}[
    cell/.style={circle, draw=#1, fill=#1!30, minimum size=0.5cm, line width=0.8pt},
    token/.style={rectangle, draw=#1, fill=#1!25, rounded corners=2pt,
                  minimum width=1.2cm, minimum height=0.6cm, font=\tiny\sffamily\bfseries},
    milblock/.style={rectangle, draw=milcolor!80!black, fill=milcolor!30, rounded corners=3pt,
                     minimum width=1.8cm, minimum height=0.8cm, font=\tiny\sffamily},
    attnblock/.style={rectangle, draw=attncolor!80!black, fill=attncolor!20, rounded corners=3pt,
                      minimum width=2cm, minimum height=0.6cm, font=\tiny\sffamily\bfseries},
    arrow/.style={->, >=Stealth, line width=0.6pt, #1}
]

% Title
\node[font=\normalsize\sffamily\bfseries] at (0, 6) {EA-MIST: Evolution-Aware Multiple Instance Set Transformer};

% ============================================
% LAYER A: SPATIAL NICHE
% ============================================
\node[font=\scriptsize\sffamily\bfseries, anchor=west] at (-6, 5) {Layer A: Spatial Niche};

% Receiver cell
\node[cell=receivercolor, minimum size=0.8cm] (recv_cell) at (-4, 4) {};
\node[font=\tiny\sffamily, below=0.05cm of recv_cell] {Receiver};

% Ring 1 cells (closest)
\foreach \i/\a in {1/30, 2/90, 3/150, 4/210, 5/270, 6/330} {
    \node[cell=ringcolor!80] (r1_\i) at ($(recv_cell) + (\a:1.2cm)$) {};
}
\node[font=\tiny\sffamily, ringcolor] at (-4, 2.3) {Ring 1 (0-50$\mu$m)};

% Ring 2 cells (farther)
\foreach \i/\a in {1/15, 2/75, 3/135, 4/195, 5/255, 6/315} {
    \node[cell=ringcolor!50] (r2_\i) at ($(recv_cell) + (\a:2cm)$) {};
}
\node[font=\tiny\sffamily, ringcolor!70] at (-4, 1.5) {Ring 2 (50-100$\mu$m)};

% ============================================
% MIL AGGREGATION (per ring)
% ============================================
\node[milblock] (mil1) at (-0.5, 4.5) {Set Transformer\\(ISAB/PMA)};
\node[milblock] (mil2) at (-0.5, 3.5) {Set Transformer\\(ISAB/PMA)};

% Arrows from rings to MIL
\draw[arrow=ringcolor, dashed] (-2.5, 4.5) -- (mil1);
\draw[arrow=ringcolor!70, dashed] (-2.5, 3.5) -- (mil2);

\node[font=\tiny\sffamily, text width=1.5cm, align=center] at (-1.5, 4.5) {$k_1$ cells};
\node[font=\tiny\sffamily, text width=1.5cm, align=center] at (-1.5, 3.5) {$k_2$ cells};

% Ring tokens output
\node[token=ringcolor] (ring1_tok) at (1.8, 4.5) {Ring 1};
\node[token=ringcolor!70] (ring2_tok) at (1.8, 3.5) {Ring 2};

\draw[arrow=ringcolor] (mil1) -- (ring1_tok);
\draw[arrow=ringcolor!70] (mil2) -- (ring2_tok);

% ============================================
% LAYER A: REFERENCE EMBEDDINGS
% ============================================
\node[font=\scriptsize\sffamily\bfseries, anchor=west] at (-6, 2.5) {Layer A: References};

% Reference embedding boxes
\node[token=hlcacolor, minimum width=1.5cm] (hlca_emb) at (-3, 1.5) {HLCA\\(30D)};
\node[token=lucacolor, minimum width=1.5cm] (luca_emb) at (-1, 1.5) {LuCA\\(10D)};

% Project to tokens
\node[token=hlcacolor] (hlca_tok) at (1.8, 1.8) {HLCA};
\node[token=lucacolor] (luca_tok) at (1.8, 1.0) {LuCA};

\draw[arrow=hlcacolor] (hlca_emb) -- (hlca_tok);
\draw[arrow=lucacolor] (luca_emb) -- (luca_tok);

\node[font=\tiny\sffamily] at (0.3, 1.8) {project};
\node[font=\tiny\sffamily] at (0.3, 1.0) {project};

% ============================================
% RECEIVER TOKEN
% ============================================
\node[token=receivercolor] (recv_tok) at (1.8, 5.5) {Receiver};
\draw[arrow=receivercolor] (recv_cell) -- ++(2, 1) -- (recv_tok);
\node[font=\tiny\sffamily] at (-1, 5.2) {encode};

% ============================================
% LAYER B+C: UNIFIED SELF-ATTENTION
% ============================================
\node[font=\scriptsize\sffamily\bfseries, anchor=west] at (3.5, 5.5) {Layer B+C: Self-Attention};

% Token sequence bracket
\draw[decorate, decoration={brace, amplitude=5pt}] (3.5, 5.8) -- (3.5, 0.7)
    node[midway, right=8pt, font=\tiny\sffamily, text width=1.2cm, align=left] {
        9-token\\sequence
    };

% Self-attention block
\node[attnblock, minimum width=2.5cm, minimum height=3cm] (sattn) at (5.5, 3.2) {};
\node[font=\tiny\sffamily\bfseries, rotate=90] at (5.5, 3.2) {Multi-Head Self-Attention};

% Connect all tokens to attention
\draw[arrow=gray] (recv_tok) -- (sattn);
\draw[arrow=gray] (ring1_tok) -- (sattn);
\draw[arrow=gray] (ring2_tok) -- (sattn);
\draw[arrow=gray] (hlca_tok) -- (sattn);
\draw[arrow=gray] (luca_tok) -- (sattn);

% ============================================
% OUTPUT
% ============================================
\node[token=receivercolor, minimum width=2cm] (output) at (8, 3.2) {Receiver\\State};
\draw[arrow=receivercolor, line width=1.2pt] (sattn) -- (output);

% ============================================
% ANNOTATIONS
% ============================================

% MIL explanation
\node[rectangle, draw=gray!50, fill=gray!5, rounded corners, font=\tiny\sffamily,
      text width=2.5cm, align=left] at (-4, -0.3) {
    \textbf{MIL Aggregation:}\\
    Permutation invariant\\
    over cells in each ring.\\
    Variable $k$ per ring.
};

% Attention explanation
\node[rectangle, draw=gray!50, fill=gray!5, rounded corners, font=\tiny\sffamily,
      text width=2.8cm, align=left] at (5.5, -0.5) {
    \textbf{Self-Attention:}\\
    All tokens attend to all.\\
    HLCA/LuCA provide\\
    healthy vs cancer bias.\\
    Learns niche influence.
};

% AMICI connection
\node[rectangle, draw=black!50, fill=yellow!10, rounded corners, font=\tiny\sffamily,
      text width=3.5cm, align=left] at (8, 0) {
    \textbf{From AMICI:}\\
    $Q$: Receiver identity\\
    $K$: Neighbor phenotypes\\
    $V$: Neighbor influence\\
    + Reference geometry
};

\end{tikzpicture}
\end{document}
